\section{Neuro-Symbolic Retrieval}
\label{sec:retrieval}
In this section, we introduce a neuro-symbolic retrieval algorithm featuring tri-path collaborative entity recall and achieve deep topological reasoning through graph traversal.
It can also serve as a fundamental retrieval algorithm adaptable to various downstream tasks in \S\ref{sec:application}.

\subsection{Node Matching}
Our retrieval system supports arbitrary query formats, including keywords, scientific questions, abstracts, idea texts, and even complete papers.
Given a query $q$, we map it into KG nodes through three distinct ways.

\paragraph{Keyword Matching.}
We use an LLM to extract keywords from $q$ and assign each keyword with an importance score, forming a keyword list $\mathcal{K}=\{(k_i, s_i^\text{llm})\}_{i=1}^m$, where $k_i$ is the $i$-th extracted keyword with text normalization and $s_i^\text{llm} \in [0, 1]$ represents its normalized importance score.
The maximum number of keywords extracted by the LLM is $m$.
Then, we first perform exact text matching of $k_i$ in the KG.
For each matched keyword node $g$, we assign it an exact match score:
\begin{align}
    \text{score}_{exact}(k_i, g)=s_i^\text{llm}.
\end{align}
Second, we perform vector matching.
After encoding each $k_i$ into a semantic vector, we compute semantic similarity based on the pre-calculated keyword text embeddings in the KG.
Nodes with similarity scores exceeding the threshold $\theta_{kw}$ (default to $0.7$) are retained, with their scores as:
\begin{align}
    \text{score}_{vec}(k_i, g)=s_i^\text{llm}\cdot\text{sim}(k_i, g).
\end{align}
If multiple nodes surpass the threshold, we select only the top-3 nodes for each $k_i$.
The same keyword node $g$ may be matched by multiple input keywords or simultaneously by both exact and vector matching.
We take the maximum of all its scores as the node's final weight:
\begin{align}
    w_g^{kw}=\max_i \left( \mathbb{1}[k_i=g]\cdot s_i^\text{llm}, \mathbb{1}[\text{sim}(k_i,g)\ge \theta_{kw}]\cdot s_i^\text{llm}\text{sim}(k_i, g)\right)
\end{align}
The final set of keyword-matching nodes is denoted as $\mathcal{K}_{seed}=\{(g, w_g^{kw})\}$.

\paragraph{Semantic Matching.}
We embed query $q$ to obtain vector $\mathbf{e}_q$ (Here, if the input is an entire paper, we only extract its abstract for embedding.), which is then used to retrieve the top-$60$ papers from the KG based on title embeddings and abstract embeddings, respectively.
We then employ a reranker (\texttt{bge-reranker-large} \citep{bge}) to re-rank the retrieved papers, retaining the top-$15$ papers for title and abstract.
Given a retrieved paper $p$, we define $s_p^{title}$ and $s_p^{abs}$ as its retrieval scores through title or abstract matching, and compute a weighted combination of the two scores:
\begin{align}
\label{eq:weighted}
    s_p^{emb}=\frac{0.4\cdot s_p^{title}+0.6\cdot s_p^{abs}}{0.4\cdot \mathbb{1}[\exists s_p^{title}]+0.6\cdot \mathbb{1}[\exists s_p^{abs}]}.
\end{align}
Here, it is set to $0$ if $s_p^{title}$ or $s_p^{abs}$ does not exist.
The final candidate paper nodes from semantic matching are denoted as $\mathcal{P}^{emb}\{(p, s_p^{emb})\}$.

\paragraph{Title Matching.}
Since titles encapsulate the most critical information of papers and are highly beneficial for paper retrieval, we specifically perform title matching for queries $q$ that contain titles.
We use GROBID to extract all titles (including the paper's title and its references' titles) from the idea or paper and employ an LLM to assign a confidence score $c_j$ to each title $t_j$.
We retain the top-$10$ titles with the highest confidence scores and normalize them (removing non-alphabetic characters and converting to lowercase) to obtain the title set $\mathcal{T}=\{(t_j, c_j)\}_{j=1}^n$.
We then perform text matching of titles in the KG.
If an exact match is found, a matching score of $m(t_j, p)=1.0$ is assigned; otherwise, we compute the fuzzy similarity between two titles based on the following formula:
\begin{align}
    m(t_j, p)=0.65\cdot \text{seq}(t_j, p)+0.35\cdot \text{token\_overlap}(t_j, p),
\end{align}
where $\text{seq}(a, b)$ is based on the Longest Common Subsequence (LCS) of $a$ and $b$, and $\text{token\_overlap}$ computes the Jaccard overlap ratio of the token sets of $a$ and $b$.
Candidates with similarity below $\theta_{title}$ (default to $0.88$) are directly discarded.
For paper $p$ matched by title $t_j$, we assign it a score:
\begin{align}
    s_{j,p}^{title}=c_j\cdot m(t_j,p).
\end{align}
If the same paper is matched by multiple titles, we take the maximum score $s_p^{title}=\max_js_{j,p}^{title}$.
Each input title retains at most the top-5 papers, and all papers constitute $\mathcal{P}^{title}=\{(p, s_p^{title})\}$.

\paragraph{Node Merging.}
We obtain two candidate paper node sets through the semantic and title pathways.
Then we need to merge them into $\mathcal{P}_{seed}$ and unify their weights.
For each candidate paper $p$, we compute the dot product with vector $\mathbf{e}_q$ and apply weighting according to the ratio specified in Eq.\ref{eq:weighted}:
\begin{align}
    \bar{s}_p^{emb} = \text{combine}(\text{sim}_p^{title},\text{sim}_p^{abs}),\quad\text{sim}_p^{title} = \mathbf{e}_q^\top \mathbf{e}_p^{title},\quad\text{sim}_p^{abs} = \mathbf{e}_q^\top \mathbf{e}_p^{abs}.
\end{align}
We then perform MinMax normalization:
\begin{align}
    \widetilde{s}_p^{emb} = \text{MinMaxNorm}(\bar{s}_p^{emb}),\quad\widetilde{s}_p^{title} = \text{MinMaxNorm}(s_p^{title}),\quad\text{MinMaxNorm}(x_p) = \frac{x_p - x_{min}}{x_{max} - x_{min}}
\end{align}
Finally, we define the unified paper weight:
\begin{align}
    \label{eq:pre_graph}s_p^{pre}=\lambda_{emb}\widetilde{s}_p^{emb}+\lambda_{title}\widetilde{s}_p^{title}+b_p^{pre},\quad b_p^{pre} =
\begin{cases}
0.35, & \text{exact title hit} \\
0.10, & \text{fuzzy title hit} \\
0, & \text{otherwise}
\end{cases},
\end{align}
where $b_p^{pre}$ denotes the title bonus, and $\lambda_{emb}$ (default to $0.3$) and $\lambda_{title}$ (default to $0.8$) represent the importance weights for semantic and title pathways, respectively.

\subsection{Weight Setting}
Taking $\mathcal{K}_{seed}$ and $\mathcal{P}_{seed}$ as starting points, we perform a 2-hop subgraph propagation, where all edges are treated as undirected during the propagation process.
To prevent subgraph explosion, we select at most $500$ nodes per hop for each entity type.
For each paper $p$ in the local subgraph, we compute its importance based on its citation count $c_p$.
Let $C$ denote total citation counts for all papers in the subgraph.
The paper's importance is defined as:
\begin{align}
\label{eq:imp}
    \text{imp}(p)=\min\left(1,\frac{\log(1+c_p)}{\log(1+\max(1,C))}\right).
\end{align}
Here, the importance can be tailored to the downstream task: if the task emphasizes paper quality, it can be computed according to Eq.\ref{eq:imp}; if the focus is solely on relevance, all papers can be forced to $\text{imp}(p)=1$.
For each seed paper $p$, we define its unnormalized weight as:
\begin{align}
\label{eq:wp}
    w_p^{seed}=s_p^{pre}\cdot(1+\gamma\cdot \text{imp}(p)),
\end{align}
where $\gamma$ is the control factor for importance (default to $0.5$).
For each seed keyword $g$, we define its unnormalized weight as $w_g^{seed} = w_g^{kw}$.
We define the distribution $\mathbf{s}$ over all nodes in the graph as:
\begin{align}
    s_v =
\begin{cases}
\dfrac{w_v^{seed}}{Z}, & v \in S \\
0, & v \notin S
\end{cases},\quad Z = \sum_{v \in S} w_v^{seed},\quad S = \mathcal{P}_{seed} \cup \mathcal{K}_{seed}.
\end{align}

\begin{table*}[!t]
\small
\centering
\caption{Definitions of Unnormalized Edge Weights ($\omega(u,v)$).}
\renewcommand{\arraystretch}{1.}
\begin{tabular}{@{} l| p{6cm} p{7cm} @{}}
\toprule
\textbf{Edge Type} & \textbf{Weight Formula(s)} & \textbf{Parameter Description} \\ 
\midrule
\texttt{HAS\_KEYWORD} & 
$\begin{aligned}[t]
&\omega_{\text{HK}}(p,g) = \beta_{hk} \cdot \kappa(g) \cdot \text{rel}_{p,g} \\
&\kappa(g) = \begin{cases} 
w_g^{kw}, &\text{if } g \text{ is a seed} \\ 
\epsilon_{kw}, &\text{otherwise} 
\end{cases}
\end{aligned}$ & 
$\beta_{hk}$: Base weight for keyword association (default $1.20$). \newline
$\text{rel}_{p,g}$: Importance score from $(p,g)$. \newline
$\kappa(g)$: Prior weight modulator for the keyword node. \newline
$w_g^{kw}$: Initial matching score for seed keywords. \newline
$\epsilon_{kw}$: Smoothing factor for non-seed keywords (default $0.25$). \\
\midrule
\texttt{CITES} & 
$\omega_{\text{CITES}}(u,v) = \beta_{cite}$ & 
$\beta_{cite}$: Base weight for paper citation relation (default $1.00$). \\
\midrule
\texttt{RELATED\_TO} & 
$\omega_{\text{RELATED}}(u,v) = \beta_{rel}$ & 
$\beta_{rel}$: Base weight for paper relatedness (default $0.90$). \\
\midrule
\texttt{AUTHORED} & 
$\omega_{\text{AUTHORED}}(u,v) = \beta_{auth}$ & 
$\beta_{auth}$: Base weight for authorship relation (default $0.80$). \\
\midrule
\texttt{COAUTHOR} & 
$\begin{aligned}[t]
&\omega_{\text{COA}}(u,v) = \beta_{coauth} \cdot \max(1, \phi(n_{uv})) \\
&\phi(n_{uv}) = \min(c_{max}, \log(1+n_{uv}))
\end{aligned}$ & 
$\beta_{coauth}$: Base weight for co-authorship (default $0.60$). \newline
$n_{uv}$: Co-authoring frequency. \newline
$\phi(\cdot)$: Frequency smoothing function. \newline
$c_{max}$: Logarithmic cap to prevent infinite weight magnification (default $2.0$). \\
\midrule
\texttt{COOCCUR} & 
$\begin{aligned}[t]
&\omega_{\text{COO}}(u,v) = \beta_{cooc} \cdot \max(1, \phi(m_{uv})) \\
&\phi(m_{uv}) = \min(c_{max}, \log(1+m_{uv}))
\end{aligned}$ & 
$\beta_{cooc}$: Base weight for keyword co-occurrence (default $0.60$). \newline
$m_{uv}$: Co-occurrence frequency. \newline
$\phi(\cdot), c_{max}$: Same smoothing function and cap definition as \texttt{COAUTHOR}. \\
\bottomrule
\end{tabular}
\label{tab:edge_weight}
\end{table*}

 For an edge $e = (u, v, r)$ in the graph, we define its unnormalized weight $\omega(u,v)$ based on the edge type, as specified in Tab.\ref{tab:edge_weight}.

\subsection{Random Walk with Restart}
To more deeply explore the topological relationships between nodes and enable deep reasoning within the graph, we perform random walks on the graph based on seed nodes and edge weights.
For any node $u$, let its neighbor set be $N(u)$. The transition probability from $u$ to its neighbor $v$ is defined as:
\begin{align}
    P(v \mid u) = \frac{\omega(u,v)}{\sum_{x \in N(u)} \omega(u,x)}.
\end{align}
Assuming the node score vector at iteration $t$ is $\mathbf{r}^{(t)}$, we initialize $\mathbf{r}^{(0)}=\mathbf{s}$.
For any node $v$, its score in the next iteration is:
\begin{align}
    r_v^{(t+1)} = \alpha s_v + (1-\alpha) \sum_{u} r_u^{(t)} P(v \mid u),
\end{align}
where $\alpha$ denotes the restart probability.
If a node $u$ has no neighbors, we preserve its own mass by directly adding $(1-\alpha)r_u^{(t)}$ back to $u$ itself.
The iteration terminates when:
\begin{align}
    \|\mathbf{r}^{(t+1)} - \mathbf{r}^{(t)}\|_1 < \varepsilon,
\end{align}
where $\varepsilon = 10^{-6}$, or when the maximum number of iterations $T_{\max} = 50$ is reached.
The final graph score of node $v$ is given by $r_v = r_v^{(t^\star)}$, where $t^\star$ denotes the stopping iteration.

\subsection{Final Ranking}
Upon completing the graph propagation, the system derives a set of global node scores $\{r_v\}_{v \in V'}$ across the local subgraph.
For the purpose of paper retrieval, we isolate the scores of paper nodes:
\begin{align}
\label{eq:spgraph}
    s_p^{graph} = r_p, \quad p \in V' \cap \texttt{Paper}
\end{align}
Crucially, this stage allows for the inclusion of newly discovered paper nodes that are not part of the initial $\mathcal{P}_{seed}$ sets but are reached during graph expansion. 
% The final ranking process integrates the topological evidence with initial semantic relevance and intrinsic paper importance through the following multi-dimensional fusion.

To account for the academic impact of candidates within the retrieved context, we re-calculate the paper importance $\text{imp}_{final}(p)$ based on the citation distribution of the final candidate set.
We utilize the logarithmic scaling defined in Eq.\ref{eq:imp}, using total citations within the current pool to ensure a robust relative metric.
To prevent the graph diffusion from over-promoting distant nodes, we introduce a graph support factor $g_p$, which acts as a gating mechanism based on the initial retrieval strength:
\begin{align}
\label{eq:gp}
    g_p = \max(0.25, \tilde{s}_p^{pre}),
\end{align}
where $\tilde{s}_p^{pre}$ is the MinMax-normalized pre-graph score $s_p^{pre}$ in Eq.\ref{eq:pre_graph}.
This ensures that while graph-discovered papers can achieve high ranks, those with zero initial semantic relevance must demonstrate exceptionally strong topological support to surpass primary candidates.
We then normalize the graph score $s_p^{graph}$ with MinMax to obtain $\tilde{s}_p^{graph}$.

The comprehensive final score $s_p^{final}$ is defined as a weighted linear combination of three normalized components and a title-matching bonus:
\begin{align}
\label{eq:final}
    s_p^{final} = \min\left(1, \lambda_{pre}\tilde{s}_p^{pre} + \lambda_{graph}\tilde{s}_p^{graph} g_p + \lambda_{imp}\text{imp}_{final}(p)\right),
\end{align}
where $\lambda_{pre}$ (default to $0.35$) is the initial relevance, $\lambda_{graph}$ (default to $0.45$) is the topological support from graph, and $\lambda_{imp}$ (default to $0.20$) is the citation importance.
We finally return the \textbf{top-$20$ papers}, accompanied by a \textbf{detailed score breakdown} and \textbf{path-based explanations} to provide researchers with a transparent and deterministic ``cognitive map'' of the retrieval results.
The entire retrieval process can be completed within 2 minutes, significantly shorter than LLM-based deep research frameworks, while still delivering high-relevance results with in-depth topological reasoning.